\section{绪论}
\subsection{研究背景}
长期资产配置的决策对象是不同风险来源之间的资本分配。资产数量增加，为分散组合提供了条件，也增加了估计与实施的复杂性。同属股票的宽基和行业工具可能在压力期同时下跌；债券资产虽然日常波动较小，仍受到利率变化影响；境内上市的海外资产基金又同时受到标的价格、汇率和基金折溢价影响。因此，配置研究需要区分证券名称、资本权重和经济风险，不能将持仓数量直接等同于分散程度。

ETF 为多资产研究提供了相对清晰的交易工具，但指数投资的便利并不消除实施条件。不同基金的上市时间、有效历史长度和交易日历不一致。若直接将目前可见的完整资产池回填至历史起点，模型便会使用当时尚不存在的工具。若为获得完整矩阵而压缩评价区间，又可能改变原本需要解释的市场经历。本文将资产池固定，将每期可用资产集合交由历史信息决定，在共同评价日期上比较最终收益路径。

传统配置分析常将收益和风险写为相互竞争的两个目标。实际研究还需要加入第三个维度，即结果能否在既定记账与交易规则下复现。高夏普可能来自很小的波动分母，而较高年化收益也可能伴随难以承受的回撤。仅凭单个指标无法识别这种差异。本文同时报告收益水平、波动、回撤、尾部损失、换手和现金持仓，并将求解诊断与日收益记录连接起来。

\subsection{研究问题与评价目标}
本文围绕 Global RRP 的指定配置展开研究。该配置采用 \lookbackDays{} 日历史窗口、收缩协方差及收益半衰期 20 日，并以年度滚动方式校准两项惩罚系数。组合保持周频、多头和无杠杆。正文只展示这一规格及四个对照模型的结果，规格的历史选择属性仍予以保留。

研究首先讨论风险预算参考及最终权重的凸优化结构，并说明收益目标与惩罚系数如何由当期可行参考和历史数据确定。随后分析不同时间上市的资产如何进入窗口，收益均值和协方差采用怎样的缺失数据规则，估计时点如何与收益记账衔接。实证部分检验指定组合的收益、风险和成本，同时核查全部合格资产是否确实得到使用。

本文以净年化收益、年化波动、夏普、Sortino、最大回撤、换手和成本共同评价历史结果。优化过程使用条件于历史数据的估计量，回测评价使用随后实现的收益，二者不能互相替代。

\subsection{研究对象与比较范围}
正式比较共五个模型。Global RRP 为主模型，HRP Benchmark、HERC Benchmark、Equal Weight 和 60/40 Benchmark 为对照。主模型每周调仓，对照按现有月频规则运行。所有模型使用相同评价区间、无风险利率和成本费率，但方法和调仓日历并未完全相同。因此，最终绩效差异应解释为完整实施规则的差异，不单独归因为某一优化项。

本文的“全球”指通过中国市场可交易工具获得多地区、多资产敞口，并非直接在全球各交易所同步执行。评价收益以仓库保存的基金复权价格为基础，未构造汇率对冲账户，也未估计不同市场收盘时差下的实际成交价格。分析范围限定为资产配置研究，不能将其扩展为已有成交验证的交易系统。

\subsection{研究路线与内容安排}
研究从可复现的数据集合出发，形成逐期合格资产集合，再估计收益和协方差，依次求解风险预算参考及最终权重。调仓后的权重进入日度记账系统，交易成本按交易前漂移权重计算。汇总指标、年度结果和持仓结构均来自同一条日收益路径，避免在不同章节中混用不同回测版本。

后续章节依次讨论相关研究、数学结构、估计与数据、执行及评价规则、五模型实证、指定组合的持仓解释和机构应用边界。附录仅提供资产池、符号和复现文件，不列逐周持仓明细。完整持仓 CSV 保存在仓库，使正文保持可读性，同时保留对单期结果进行核对的条件。

\subsection{研究贡献与限制}
本文的贡献主要体现在方法与实现的对应。凸风险预算参考为最终配置提供可解释的起点，正定化后的协方差支持稳定求解，收益目标直接取参考组合的预测收益，从而获得当期可行基准。两项惩罚按严格早于生效日的历史区间更新，资本约束、收益估计和成本记账分别承担明确作用。
